\section{容错与诊断}
在设计该系统时，我们面临的最大挑战之一，是处理频繁发生的组件故障。
组件的质量和数量共同决定了：这些问题是常态而非例外。我们既不能完全
信任机器，也不能完全信任磁盘。
组件故障可能导致系统不可用，或者更糟的是导致数据损坏。下面将讨论我们
如何应对这些挑战，以及当问题不可避免地发生时，系统内置了哪些用于诊断
问题的工具。

\subsection{高可用性}
在GFS集群的数百台服务器中，任意时刻都必然有部分服务器不可用。我们
通过两种简单但有效的策略保持整个系统的高可用性：快速恢复和复制。

\subsubsection{快恢复}
无论是master还是chunkserver，都被设计为无论以何种方式终止，都能在数秒
内恢复状态并启动。事实上，我们不区分正常终止和异常终止；服务器通常只需
直接杀死进程即可关闭。
客户端和其他服务器只会经历一次轻微停顿：它们等待尚未完成的请求超时，
重新连接到已重启的服务器，然后重试。section6.2.2报告了实际观察到的
启动时间。

\subsubsection{chunk复制}
如前所述，每个chunk都会复制到位于不同机架的多个chunkserver上。用户
可以为文件命名空间的不同部分指定不同的复制级别，默认值为3。
当chunkserver离线，或通过校验和验证检测到replica损坏时(见section5.2)，
master会按需克隆现有replica，以确保每个chunk始终具有完整replica。
尽管复制一直很好地满足了我们的需求，但随着只读存储需求不断增长，我们
正在探索其他跨服务器冗余形式，例如奇偶校验或纠删码。
由于我们的流量主要由追加和读取构成，而不是小型随机写入，我们预计在
这个非常松耦合的系统中实现这些更复杂的冗余方案虽有挑战，但仍然可控。

\subsubsection{master复制}
master状态通过复制来保证可靠性。它的操作日志和检查点会复制到多台机器
上。只有当状态mutation对应的日志记录已在本地和所有master replica上刷新到
磁盘后，该mutation才被视为已提交。
为简单起见，一个master进程始终负责全部mutation，以及垃圾回收等会在
系统内部改变状态的后台活动。当它发生故障时，可以几乎立即重启。如果它
所在的机器或磁盘发生故障，GFS外部的监控基础设施会在其他位置利用已复制
的操作日志启动一个新的master进程。
客户端只使用master的规范名称(例如gfs-test)。该名称是一个DNS别名；
如果master被迁移到另一台机器，该别名可以被修改。

此外，即使primary master宕机，“shadow”master仍可向文件系统提供只读
访问。它们是shadow而非mirror，因为它们可能略微落后于primary，通常只
落后几分之一秒。
对于未被积极mutation的文件，或不介意得到略旧结果的应用程序，shadow
master能够提高读取可用性。事实上，由于文件内容是从chunkserver读取的，
应用程序不会观察到过期的文件内容。
在短暂窗口中可能过期的是文件元数据，例如目录内容或访问控制信息。

为了保持信息更新，shadow master会读取持续增长的操作日志replica，并像
primary一样，将完全相同的变更序列应用到自己的数据结构中。与primary
相同，它会在启动时(以及此后不频繁地)轮询chunkserver以定位chunk replica，
并与它们频繁交换握手消息来监控状态。
shadow master只依赖primary master提供replica位置更新；这些更新源于
primary关于创建和删除replica的决策。

\subsection{数据完整性}
每个chunkserver都使用校验和来检测已存储数据的损坏。考虑到一个GFS集群
通常在数百台机器上拥有数千块磁盘，它会经常遇到磁盘故障；这些故障会在
读取和写入路径上造成数据损坏或丢失。(其中一个原因见section7。)
我们可以使用其他chunk replica从损坏中恢复，但通过比较不同chunkserver上的
replica来检测损坏并不现实。此外，replica之间存在差异也可能是合法的：GFS的
mutation语义，特别是前文讨论的原子\textbf{record append}，并不保证replica
完全相同。因此，每个chunkserver都必须通过维护校验和，独立验证自身replica
的完整性。

一个chunk被划分为64KB的块，每个块都有对应的32位校验和。和其他元数据
一样，校验和保存在内存中，并通过日志持久化存储；它们与用户数据分开保存。

对于读取操作，chunkserver会在向请求方返回任何数据前，验证与读取范围
重叠的数据块的校验和。无论请求方是客户端还是另一个chunkserver，都遵循
这一规则。因此，chunkserver不会将损坏传播到其他机器。
如果某个块与记录的校验和不匹配，chunkserver会向请求方返回错误，并向
master报告该不匹配。随后，请求方会从其他replica读取，master则会从另一个
replica克隆该chunk。新的有效replica就位后，master会指示报告不匹配的
chunkserver删除自己的replica。

出于几个原因，校验和对读取性能影响很小。由于大多数读取至少跨越数个块，
为了验证，我们只需额外读取并计算校验和的相对少量数据。GFS客户端代码还
会尽量将读取对齐到校验和块边界，以进一步降低这部分开销。
此外，chunkserver上的校验和查找与比较不需要任何I/O，而校验和计算通常
能够与I/O操作重叠执行。

对于追加到chunk末尾的写入(而非覆盖已有数据的写入)，校验和计算得到了
高度优化，因为它们在我们的工作负载中占主导地位。我们只需增量更新最后
一个未填满的校验和块的校验和，并为追加操作填满的任何全新校验和块计算
新的校验和。
即使最后一个未填满的校验和块已经损坏，而我们此时未能检测到，新校验和
值也不会与已存储数据匹配；下次读取该块时，仍会像往常一样检测到损坏。

相反，如果写入覆盖了chunk中的一个现有范围，我们必须先读取并验证被
覆盖范围的第一个和最后一个块，再执行写入，最后计算并记录新的校验和。
如果在部分覆盖首尾块前不验证它们，新校验和可能会掩盖未被覆盖区域中
原本存在的损坏。

在空闲期间，chunkserver可以扫描并验证不活跃chunk的内容。这使我们能够
检测很少被读取的chunk中的损坏。一旦检测到损坏，master就可以创建一个
新的未损坏replica，并删除损坏replica。
这样可以避免一个不活跃但已损坏的chunk replica误导master，使其以为某个
chunk已经拥有足够数量的有效replica。

\subsection{诊断工具}

广泛且细致的诊断日志，对问题隔离、调试和性能分析带来了极大帮助，
而成本却很低。没有日志，很难理解机器之间短暂且无法复现的交互。
GFS服务器会生成诊断日志，记录许多重要事件(例如chunkserver上线和
下线)，以及全部RPC请求和回复。删除这些诊断日志不会影响系统正确性。
不过，只要存储空间允许，我们会尽量保留这些日志。

RPC日志包含网络上传输的精确请求与响应，但不包含正在读取或写入的文件
数据。通过匹配请求和回复，并整理不同机器上的RPC记录，我们可以重建
完整的交互历史，从而诊断问题。
这些日志也可作为负载测试和性能分析的追踪记录。

日志记录对性能的影响很小(且远小于它带来的收益)，因为这些日志以顺序、
异步方式写入。最近发生的事件还会保存在内存中，可供持续在线监控。